\scanpage{54}
% letterhead
\begin{center}
{\sffamily HARVARD UNIVERSITY\\
CAMBRIDGE, MASSACHUSETTS 02138}
\end{center}
% upper-right pencil note, with its open, irregular frame
\margin{\figure{\begin{tikzpicture}[baseline=(copie.base)]
\node[align=left,inner sep=2pt] (copie) {Faire copie\\pour \unsure{Jouanolou}};
\draw[thin] ([xshift=-3pt,yshift=1pt]copie.north west)
 .. controls ([yshift=3pt]copie.north) and ([xshift=-2pt,yshift=5pt]copie.north east)
 .. ([xshift=-2pt,yshift=2pt]copie.north east)
 -- ([xshift=-2pt,yshift=-3pt]copie.north east);
\draw[thin] ([yshift=-1pt]copie.north west) -- (copie.south west)
 -- ([xshift=3pt,yshift=1pt]copie.south east)
 -- ([xshift=4pt,yshift=-7pt]copie.north east);
\end{tikzpicture}}}
% letterhead
\noindent {\sffamily DEPARTMENT OF MATHEMATICS\\
2 DIVINITY AVENUE}\\
\hfill Aug 4, 1969\\

\noindent Dear Grothendieck,\\

I have been reading recently some notes from the IHES (one of your seminars, I suppose) entitled ``Exposé \(\overline{\underline{\mathrm{V}}}\). Systèmes
\reversemarginpar\margin{\figure{\begin{tikzpicture}\draw[thin] (0,0) -- (0,-1.15);\end{tikzpicture}}}\normalmarginpar%
projectifs J-adiques, par J.P. Jouanolou.'' Could you tell me (for purposes of reference) to which seminar they belong, and whether there is a more recent edition?

Also I have some comments and questions. The most serious question is that I do not understand the proof of the ``Theorem of Shih'', \unsure{A.3.1|A.3.1.} on page 46. It seems to me that the lemma on p. 45 is valid only for \(r>p\) (cf EGA \(\unsure{0|O}_{\overline{\underline{\mathrm{III}}}}\) 13.5.5.2), and so \deleted{\unsure{the|from}} the equivalence of statements (i) and (ii) of the theorem is not clear. (I do agree, however, with the weaker statement ``(resp.\ldots)'' in parentheses). This casts doubt on the following theorem A.3.2 on p. 47, and also on the earlier Prop. 5.3.1 on p. 40\unsure{.|,} These difficulties would disappear if one worked \unsure{throughout} in the category \(\mathrm{pro}\text{-}A\)
% Two tiny, obliquely aligned marginal points; punctuation function uncertain.
\margin{\figure{\tikz{\node[rotate=-15,inner sep=0pt]{\unsure{:}};}}}
of projective systems with the usual equivalence, rather than in the \deleted{category} finer category \(P_{AR}\).

Other comments: p. 42 line \(\dot{-}1\). Isn't the spectral sequence actually biregular in case of a finite filtration?

p. 47 line \(-1\). I think the hypothesis should be stated for \(n\) and \(n+1\), instead of \(n\) and \(n-1\) (cf EGA \(0_{\overline{\underline{\mathrm{III}}}}\) 11.1.10).
% archive sheet number: 53
\clearpage
